\chapter{Diseño, cálculos y simulaciones}
\label{chap:diseno-calculos}
La señal de entrada se puede ver en la figura \ref{plot:vi}. Cabe destacar que el análisis de la simulación comienza a
los $9.05 \si{\milli\second}$ para permitir que se estabilicen todos los elementos. Este arranque desfasado se mantiene
para todos los gráficos de la simulación.

\begin{figure}[!ht]
  \centering
  \input{figures/plot_vi.tex}
  \caption{señal de entrada.}
  \label{plot:vi}
\end{figure}

\section{Etapa 1: Adaptación de impedancia}
La impedancia de entrada debe ser de $\SI{15}{\kilo\ohm}$. Para ello, se utilizó la configuración buffer o seguidor de
tensión de los amplificadores operacionales. Debido a que la impedancia de entrada está por el orden de los megaohm,
podemos decir que es despreciable y colocar cualquier resistencia en paralelo entre la entrada y la masa para generar la
adaptación de impedancia deseada. El circuito a implementar para esta etapa se puede ver en la figura
\ref{crkt:adaptacion.impedancia}.

\begin{figure}[!ht]
  \centering
  \begin{tikzpicture}[transform shape]
	% Paths, nodes and wires:
	\node[op amp, yscale=-1] at (0.198, -0.01){};
	\draw (-2.25, -2) -| (-2.25, -0.5) -- (-1, -0.5);
	\draw (2, -2) |- (1.388, -0.01);
	\draw (-2.25, -2) -- (2, -2);
	\draw (-0.992, 0.48) -| (-5.25, 0.5);
	\draw (-4.25, 0.5) to[american resistor, l={$15 \si{\kilo \ohm}$}] (-4.25, -2);
	\node[ground] at (-4.25, -2){};
	\node[ocirc](N1) at (-5.25, 0.5){} node[anchor=south] at (N1.north){$V_i$};
	\draw (2, -0) -- (3.5, -0);
	\node[ocirc](N2) at (3.5, -0){} node[anchor=south] at (N2.north){$V_{buffered}$};
\end{tikzpicture}
  \caption{circuito amplificador operacional en configuración seguidor con impedancia de entrada especificada.}
  \label{crkt:adaptacion.impedancia}
\end{figure}

La señal resultante de esta etapa es idéntica a la señal de entrada.

\section{Etapa 2: Ajuste de la tensión de offset}
La señal de salida de la 2da etapa debe ser igual a la de la entrada con un offset de $-1 \si{\volt}$. Para lograr
esto, se usa el amplificador no inversor con una tensión de referencia. Esta tensión debería ser obtenida mediante un
dispositivo que provea una referencia constante, como un diodo zenner programable (ej.: TL431). Debido a que este
trabajo practico no requiere tanta exactitud en los valores de tensión, optamos por usar divisores resistivos como
referencias. Como el amplificador operacional no genera una carga significativa (debido a la altísima impedancia de
entrada) sobre el divisor resistivo, no hay problema en usarlo. En la implementación se utilizara un
potenciómetro en estos divisores resistivos para poder mitigar las tolerancias de fabricación de las resistencias mismas
y cualquier resistencia de contacto. El circuito a implementar para esta etapa se puede ver en la figura
\ref{crkt:offset}.

\begin{figure}[!ht]
  \centering
  \begin{tikzpicture}[transform shape]
	% Paths, nodes and wires:
	\node[op amp, yscale=-1] at (0.198, -0.01){};
	\draw (-2.25, -1.75) -| (-2.25, -0.5) -- (-1, -0.5);
	\draw (2, -1.75) |- (1.388, -0.01);
	\node[ocirc](N1) at (-6, 0.48){} node[anchor=south] at (N1.north){$V_{buffered}$};
	\draw (2, -0.01) -| (3.5, -0);
	\node[ocirc](N2) at (3.5, -0){} node[anchor=south] at (N2.north){$V_{offseted}$};
	\draw (-5, 0.48) to[american resistor, l={$100k\Omega$}] (-3, 0.48);
	\draw (-3, 0.48) to[american resistor, l_={$100k\Omega$}] (-3, 2.48);
	\draw (-2.25, -0.5) to[american resistor, l_={$100k\Omega$}] (-3.75, -0.5);
	\draw (-4, -1.25) to[american resistor, l_={$10k\Omega$}] (-6, -1.25);
	\draw (-1, -1.75) to[american resistor, l={$100k\Omega$}] (1, -1.75);
	\draw (-3.75, -1.25) to[american potentiometer, l={$\approx909.091\Omega$}, label distance=-0.02cm, v_=$1V$, voltage/distance from node=0.7, voltage/bump b=4.2, voltage/shift=1.4, voltage=american] (-3.75, -3);
	\draw (-3.19, -2.125) |- (-3.75, -1.25) -- (-4, -1.25);
	\draw (-3.75, -0.5) -| (-3.75, -1.25);
	\draw (-2.25, -1.75) -- (-1, -1.75);
	\draw (1, -1.75) -- (2, -1.75);
	\draw (-5, 0.48) -- (-6, 0.48);
	\node[ground] at (-3.75, -3){};
	\node[ground] at (-3.75, -3){};
	\node[ocirc](N3) at (-6.5, -1.25){} node[anchor=south] at (N3.north){$12V$};
	\node[ground, xscale=-1, yscale=-1] at (-3, 2.48){};
	\draw (-0.992, 0.48) -- (-3, 0.48);
	\draw (-6, -1.25) -- (-6.5, -1.25);
\end{tikzpicture}
  \caption{circuito amplificador operacional en configuración amplificador no inversor con tensión de offset.}
  \label{crkt:offset}
\end{figure}

La simulación de esta etapa dio como resultado la señal que se puede ver en la figura \ref{plot:offset}.

\begin{figure}[!ht]
  \centering
  \input{figures/plot_offset.tex}
  \caption{señal de entrada con offset.}
  \label{plot:offset}
\end{figure}

\section{Etapa 3: Señal diente de sierra}
Sabiendo que la señal anterior es una onda cuadrada con offset, debido a que el área de la misma crece y decrece con el
tiempo gracias al offset agregado anteriormente, podemos generar la señal diente de sierra obteniendo una tensión que
varíe a la par del área debajo de la curva, en este caso, la señal cuadrada con offset. Para ello, utilizaremos el
circuito propuesto por el profesor que se puede ver en la figura \ref{crkt:integrador}, dado que se comparten las
condiciones de operación.

\begin{figure}[!ht]
  \centering
  \begin{tikzpicture}[transform shape]
	% Paths, nodes and wires:
	\node[op amp, yscale=-1] at (-0, 0){};
	\draw (1.802, -1.74) |- (1.19, 0);
	\node[ocirc](N1) at (-4.75, -0.5){} node[anchor=south] at (N1.north){$V_{offseted}$};
	\draw (1.802, 0) -| (3.302, 0.01);
	\node[ocirc](N2) at (3.302, 0.01){} node[anchor=south] at (N2.north){$V_{integrated}$};
	\draw (-1.198, -1.74) to[american resistor, l={$100k\Omega$}] (1.802, -1.74);
	\draw (-2.19, 0.49) -- (-1.19, 0.49);
	\draw (-1.198, -1.74) -- (-1.198, -0.49);
	\draw (-1.198, -3) to[capacitor, l={$10nF$}] (1.75, -3);
	\draw (1.75, -1.75) -| (1.802, -1.74) |- (1.75, -3);
	\draw (-1.198, -3) -- (-1.198, -1.74);
	\draw (-3.5, -0.5) to[american resistor, l={$10k\Omega$}] (-1.19, -0.49);
	\draw (-4.75, -0.5) -- (-3.5, -0.5);
	\node[ground, rotate=-90] at (-2.19, 0.49){};
\end{tikzpicture}

  \caption{circuito amplificador operacional en configuración integrador.}
  \label{crkt:integrador}
\end{figure}

La simulación de esta etapa dio como resultado la señal que se puede ver en la figura \ref{plot:integrador.pre}.

\begin{figure}[!ht]
  \centering
  \pgfplotsset{
    table/col sep=semicolon,
    /pgf/number format/read comma as period
}

\begin{tikzpicture}
    \begin{axis}[
        width=.9\textwidth,
        height=5cm, % Reducido un poco más para seguridad
        xlabel={Tiempo [ms]},
        grid=both,
        xmin=0.00905, xmax=0.00952,
        scaled x ticks=false,
        xticklabel={\pgfmathparse{\tick*1000}\pgfmathprintnumber{\pgfmathresult}},
        x tick label style={/pgf/number format/fixed, /pgf/number format/precision=2},
        extra y ticks={\pgfkeysvalueof{/pgfplots/ymin},\pgfkeysvalueof{/pgfplots/ymax}},
        ylabel={Amplitud [V]}
    ]
        \addplot[blue, thick] table [x=time, y={V(/integrated_pre)}] {sources/plots/sim.csv};
    \end{axis}
\end{tikzpicture}

  \caption{señal integrada.}
  \label{plot:integrador.pre}
\end{figure}

Como la señal no cumple con los mínimos y máximos establecidos anteriormente, es necesario amplificar y agregar un
offset a la señal integrada. Para ello, se puede utilizar un amplificador operacional en configuración diferenciador,
como se puede ver en la figura \ref{crkt:diferenciador}. La ecuación \ref{eq:diferenciador.vo} describe como va a ser la
salida en función de los potenciales en la entrada inversora y no inversora.

\begin{figure}[!ht]
  \centering
  \begin{subfigure}[c]{0.6\textwidth}
    \begin{tikzpicture}[transform shape]
	% Paths, nodes and wires:
	\node[op amp, yscale=-1] at (2.5, 0.25){};
	\draw (4.302, -1.49) |- (3.69, 0.25);
	\draw (4.302, 0.25) -| (5.802, 0.26);
	\node[ocirc](N1) at (5.802, 0.26){} node[anchor=south] at (N1.north){$v_o$};
	\draw (1.302, -1.49) to[american resistor, l={$R_2$}] (3.302, -1.49);
	\draw (3.302, -1.49) -- (4.302, -1.49);
	\draw (-0.698, -0.24) to[american resistor, l={$R_1$}] (1.302, -0.24);
	\draw (-1.69, 0.74) to[american resistor, l={$R_1$}] (0.31, 0.74);
	\draw (0.31, 2.74) to[american resistor, l={$R_2$}] (0.31, 0.74);
	\draw (0.31, 0.74) -- (1.31, 0.74);
	\node[ground, xscale=-1, yscale=-1] at (0.31, 2.74){};
	\draw (1.302, -1.49) -- (1.302, -0.24);
	\draw (-1.69, 0.74) -| (-3.25, 0.75);
	\draw (-0.698, -0.24) -| (-2, -0.25);
	\node[ground] at (-2, -1.75){};
	\node[ground] at (-3.25, -1.75){};
	\draw (-2, -0.25) to[open, o-o, v=$v_1$, voltage=straight] (-2, -1.75);
	\draw (-3.25, 0.74) to[open, o-o, v=$v_2$, voltage/bump b=1.6, voltage=straight] (-3.25, -1.75);
\end{tikzpicture}
    \caption{circuito.}
    \label{crkt:diferenciador}
  \end{subfigure}
  \hfill
  \begin{subfigure}[c]{0.3\textwidth}
    \begin{equation}
      v_o = \left(v_2 - v_1 \right) \frac{R_2}{R_1}
      \label{eq:diferenciador.vo}
    \end{equation}
    \caption{ecuación.}
  \end{subfigure}
  \caption{amplificador operacional en configuración diferenciador.}
\end{figure}

Para poder obtener los potenciales y ganancias, es mas fácil ver la ecuación \ref{eq:diferenciador.vo} como dos
ecuaciones diferentes, ignorando los valores resistivos y reemplazándolos por el termino de ganancia $A_v$. De esta
forma, podemos interpretar las dos ecuaciones como un "mapeo", donde mapeamos un punto ($v_1$ o $v_2$) en otro punto
($v_o$).
\begin{equation}
  v_{o1} = \left( v_{21} - v_{11} \right) A_v \qquad v_{o2} = \left( v_{22} - v_{12} \right) A_v
\end{equation}

Dividiendo ambas ecuaciones y definiendo la constante $k = \frac{v_{o2}}{v_{o1}}$, se obtiene la relación base:
\begin{equation}
  k \cdot (v_{21} - v_{11}) = v_{22} - v_{12}
  \label{eq:base_desarrollo}
\end{equation}

Para resolver esta ecuación, definimos $S$ como los valores de la señal de entrada y $X$ como la tensión de referencia
constante. Se presentan dos configuraciones posibles:
\begin{itemize}
    \item \textbf{Caso 1 (Señal en entrada no inversora):} $v_{2x} = S_x$ y $v_{1x} = X$.
    \item \textbf{Caso 2 (Señal en entrada inversora):} $v_{1x} = S_x$ y $v_{2x} = X$.
\end{itemize}

Sustituyendo estos valores en la ecuación \ref{eq:base_desarrollo}, se observa que ambos casos conducen a la misma
estructura algebraica:
\begin{align*}
    \text{Caso 1: } & k(S_1 - X) = S_2 - X \implies kS_1 - S_2 = X(k - 1) \\
    \text{Caso 2: } & k(X - S_1) = X - S_2 \implies kS_1 - S_2 = X(k - 1)
\end{align*}

De esta forma, se llega a una única expresión para la tensión de referencia, independiente de la configuración elegida:
\begin{equation}
    X = \frac{k \cdot S_1 - S_2}{k - 1}
    \label{eq:diferenciador.offset}
\end{equation}

Finalmente, la ganancia se calcula despejándola de cualquiera de las ecuaciones iniciales: 
\begin{equation}
  A_v = \left| \frac{v_{o1}}{S_1 - X} \right|
  \label{eq:diferenciador.ganancia}
\end{equation}

El signo de la ecuación \ref{eq:diferenciador.pend.curva} determinará si la señal debe inyectarse por la entrada no
inversora ($G > 0$) o por la inversora ($G < 0$).
\begin{equation}
  G = \frac{v_{o2} - v_{o1}}{S_2 - S_1}
  \label{eq:diferenciador.pend.curva}
\end{equation}

Con las herramientas planteadas anteriormente, podemos ahora calcular la referencia y ganancias para adaptar la salida
del amplificador integrador. De antemano podemos ver que tenemos que generar una inversión de fase para hacer que el
área bajo la curva de la señal de entrada se correlacione con la tensión de salida del integrador. Para ello, tomamos el
caso 2 mencionado anteriormente, y corroboramos si realmente la señal tiene que ir en la entrada inversora utilizando la
ecuación \ref{eq:diferenciador.pend.curva}. Considerando:
\begin{align*}
  S_1 = v_{iM} = -1.81 \si{\volt} &\to v_{o1} = v_{om} = -1.5 \si{\volt} \\
  S_2 = v_{im} = -3.25 \si{\volt} &\to v_{o2} = v_{oM} = 4 \si{\volt}
\end{align*}
entonces:
\begin{align*}
  G &= \frac{4 \si{\volt} - (-1.5\si{\volt})}{-3.25\si{\volt} - (-1.81\si{\volt})} \\
  G &= -3.8 \quad \to \text{ señal en entrada inversora}
\end{align*}

Habiendo corroborado que la señal de entrada tiene que ir a la entrada inversora, podemos calcular el offset a calcular
a partir de la ecuación \ref{eq:diferenciador.offset}:
\begin{figure}[!ht]
  \centering
  \begin{subfigure}[t]{0.45\textwidth}
    \begin{align*}
      k &= \frac{v_{o2}}{v_{o1}} \\
      k &= \frac{4 \si{\volt}}{-1.5 \si{\volt}} \\
      k &= -\frac{8}{3}
    \end{align*}
  \end{subfigure}
  \hspace{.5cm}
  \begin{subfigure}[t]{0.45\textwidth}
    \begin{align*}
      X &= \frac{-\frac{8}{3} \cdot (-1.81 \si{\volt}) - (-3.25 \si{\volt})}{-\frac{8}{3} - 1} \\
      X &\approx -2.2 \si{\volt}
    \end{align*}
  \end{subfigure}
\end{figure}

En este caso, el offset es negativo. Esto significaría un problema si la fuente de alimentación del circuito
diferenciador no fuese partida, pero en este caso, como la alimentación es $\pm 12 \si{\volt}$ no hay problema, simplemente el
divisor resistivo debe ser conectado a $-12 \si{\volt}$.

Por ultimo, la ganancia la podemos calcular utilizando la ecuación \ref{eq:diferenciador.ganancia}:
\begin{align*}
  A_v &= \left| \frac{-1.5 \si{\volt}}{-1.8 \si{\volt} - (-2.2 \si{\volt})} \right| \\
  A_v &\approx 3.9
\end{align*}

El juego de resistencias que cumplen esta relación son:

\begin{equation*}
  R_1 = 100 \si{\kilo \ohm} \quad R_2 = 390 \si{\kilo \ohm}
\end{equation*}

Colocando la señal de salida del integrador en la entrada inversora y la tensión de referencia negativa en la entrada no
inversora, la simulación dio los resultados que se pueden ver en la figura \ref{plot:integrador}.

\begin{figure}[!ht]
  \centering
  \input{figures/plot_integrated.tex}
  \caption{señal integrada con amplitud y desfase correcto.}
  \label{plot:integrador}
\end{figure}

\section{Etapa 4: Volver a la señal de entrada}
Con la señal diente de sierra siendo generada a partir de la integral de la señal de entrada cuadrada, ahora debemos
volver a la señal de onda cuadrada a partir de la diente de sierra. Para "deshacer" los efectos matemáticos de la señal
diente de sierra, podemos colocar un amplificador operacional en configuración derivadora, para volver a la señal
cuadrada original. Nuevamente, para este circuito, utilizamos el propuesto por el profesor que se puede ver en la
figura \ref{crkt:derivador}.

\begin{figure}[!ht]
  \centering
  \begin{tikzpicture}[transform shape]
	% Paths, nodes and wires:
	\node[op amp, yscale=-1] at (-0, 0){};
	\draw (1.802, -1.74) |- (1.19, 0);
	\node[ocirc](N1) at (-7, -0.5){} node[anchor=south] at (N1.north){$V_{integrated}$};
	\draw (1.802, 0) -| (3.302, 0.01);
	\node[ocirc](N2) at (3.302, 0.01){} node[anchor=south] at (N2.north){$V_{derivated}$};
	\draw (-1.198, -1.74) to[american resistor, l_={$63k\Omega$}] (1.802, -1.74);
	\draw (-2.19, 0.49) -- (-1.19, 0.49);
	\draw (-1.198, -1.74) -- (-1.198, -0.49);
	\draw (-1.198, -3) to[capacitor, l_={$56pF$}] (1.75, -3);
	\draw (1.75, -1.75) -| (1.802, -1.74) |- (1.75, -3);
	\draw (-1.198, -3) -- (-1.198, -1.74);
	\draw (-5.75, -0.5) to[american resistor, l_={$4.7k\Omega$}] (-3.44, -0.49);
	\draw (-7, -0.5) -- (-5.75, -0.5);
	\node[ground, rotate=-90] at (-2.19, 0.49){};
	\draw (-3.44, -0.49) to[capacitor, l_={$1nF$}] (-1.19, -0.49);
\end{tikzpicture}
  \caption{circuito amplificador operacional en configuración derivador.}
  \label{crkt:derivador}
\end{figure}

La simulación de esta etapa dio como resultado la señal que se puede ver en la figura \ref{plot:derivated_pre}.

\begin{figure}[!ht]
  \centering
  \input{figures/plot_derivated_pre.tex}
  \caption{señal de entrada recompuesta.}
  \label{plot:derivated_pre}
\end{figure}

Nuevamente, esta señal de salida se encuentra fuera del rango solicitado, por lo que volvemos a hacer uso de las
ecuaciones \ref{eq:diferenciador.pend.curva}, \ref{eq:diferenciador.offset} y \ref{eq:diferenciador.ganancia}.
Considerando:
\begin{align*}
  S_1 = v_{im} = -3.16 \si{\volt} &\to v_{o1} = v_{oM} = 3 \si{\volt} \\
  S_2 = v_{iM} = 3.16 \si{\volt} &\to v_{o2} = v_{om} = 0 \si{\volt}
\end{align*}
entonces:
\begin{align*}
  G &= \frac{0 \si{\volt} - 3\si{\volt}}{3.16\si{\volt} - (-3.16\si{\volt})} \\
  G &= -0.47 \quad \to \text{ señal en entrada inversora}
\end{align*}

Habiendo corroborado que la señal de entrada tiene que ir a la entrada inversora, podemos calcular el offset a calcular
a partir de la ecuación \ref{eq:diferenciador.offset}:
\begin{figure}[!ht]
  \centering
  \begin{subfigure}[t]{0.45\textwidth}
    \begin{align*}
      k &= \frac{v_{o2}}{v_{o1}} \\
      k &= \frac{0 \si{\volt}}{3 \si{\volt}} \\
      k &= 0
    \end{align*}
  \end{subfigure}
  \hspace{.5cm}
  \begin{subfigure}[t]{0.45\textwidth}
    \begin{align*}
      X &= \frac{0 \si{\volt} - 3.16 \si{\volt}}{- 1} \\
      X &= 3.16 \si{\volt}
    \end{align*}
  \end{subfigure}
\end{figure}

Por ultimo, la ganancia la podemos calcular utilizando la ecuación \ref{eq:diferenciador.ganancia}:
\begin{align*}
  A_v &= \left| \frac{3 \si{\volt}}{-3.16 \si{\volt} - 3.16 \si{\volt})} \right| \\
  A_v &\approx 0.48
\end{align*}

El juego de resistencias que cumplen esta relación son:

\begin{equation*}
  R_1 = 68 \si{\kilo \ohm} \quad R_2 = 33 \si{\kilo \ohm}
\end{equation*}

Colocando la señal de salida del integrador en la entrada inversora y la tensión de referencia negativa en la entrada no
inversora, la simulación dio los resultados que se pueden ver en la figura \ref{plot:derivador}.

\begin{figure}[!ht]
  \centering
  \input{figures/plot_derivated.tex}
  \caption{señal derivada con amplitud y desfase correcto.}
  \label{plot:derivador}
\end{figure}

\section{Etapa 5: Detección de flancos de subida y bajada}
Con la señal cuadrada ya recompuesta, el siguiente paso es obtener los picos de pendiente generados en los flancos de la
señal cuadrada. Idealmente, según la definición de una onda cuadrada, la derivada de la misma seria el impulso delta
($\delta_{(t)}$). Como no es perfecta la recomposición de la señal cuadrada, la señal resultante será parecida a un pico
delta ideal. Nuevamente se utilizo el circuito propuesto por el profesor que se puede ver en la figura
\ref{crkt:2do.derivador}.

\begin{figure}[!ht]
  \centering
  \begin{tikzpicture}[transform shape]
	% Paths, nodes and wires:
	\node[op amp, yscale=-1] at (-0, -0){};
	\draw (1.75, -1.74) |- (1.138, -0.01);
	\node[ocirc](N1) at (-6.5, 0.5){} node[anchor=south] at (N1.north){$V_{derivated}$};
	\draw (1.75, -0) -| (3.25, 0.01);
	\node[ocirc](N2) at (3.302, 0.01){} node[anchor=south] at (N2.north){$V_{2nd-derivated}$};
	\draw (-1.198, -1.74) to[american resistor, l={$42k\Omega$}] (0.802, -1.74);
	\draw (0.802, -1.74) -| (1.75, -1.75);
	\node[ground] at (-3.25, -3.25){};
	\node[ground] at (-3.25, -3.25){};
	\node[ocirc](N3) at (-6.5, -1.5){} node[anchor=south] at (N3.north){$+12V$};
	\draw (-6, -1.5) -- (-6.5, -1.5);
	\draw (-3.198, -0.49) to[american resistor, l={$56k\Omega$}] (-1.198, -0.49);
	\draw (-3.25, -1.5) to[american potentiometer, l={$600\Omega$}, v=$\approx0.675V$, voltage/distance from node=1, voltage/shift=1.3, voltage=american] (-3.25, -3.25);
	\draw (-6, -1.5) to[american resistor, l={$10k\Omega$}] (-4, -1.5);
	\draw (-2.19, 0.49) -- (-1.19, 0.49);
	\draw (-1.198, -1.74) -- (-1.198, -0.49);
	\draw (-2.69, -2.375) |- (-3.242, -1.5);
	\draw (-4, -1.5) -- (-3.25, -1.5);
	\draw (-3.242, -1.5) |- (-3.198, -0.49);
	\draw (-6.5, 0.49) -- (-2.25, 0.49) -| (-2.19, 0.48);
	\draw (1.75, -1.75) to[capacitor, l={$10nF$}] (1.75, -3);
	\draw (1.75, -3) to[american resistor, l={$470\Omega$}] (1.75, -5);
	\node[ground] at (1.75, -5){};
\end{tikzpicture}
  \caption{circuito amplificador operacional en configuración derivador.}
  \label{crkt:2do.derivador}
\end{figure}

La simulación de esta etapa dio como resultado la señal que se puede ver en la figura \ref{plot:2do.derivador}.

\begin{figure}[!ht]
  \centering
  \input{figures/plot_2nd_derivated.tex}
  \caption{señal cuadrada con doble derivada.}
  \label{plot:2do.derivador}
\end{figure}

\section{Etapa 6: Rectificación}
Por ultimo, la señal solicitada debe ser la segunda derivada pero sin picos negativos, es decir, rectificada. La
diferencia está en que esta rectificación no sera hecha con elementos pasivos, si no que también con amplificadores
operacionales. Esto tiene sus ventajas y desventajas. La ventaja es que se omite el potencial de polarización requerido
por los diodos de silicio o germanio, debido a la baja tensión de offset del amplificador seleccionado. La desventaja,
es solo una señal y no es posible demandarle mucha corriente, debido a que el amplificador operacional no está preparado
para ello; debería antes pasar por otro circuito (como un amplificador) para dotar esa señal de potencia. Para realizar la
rectificación, se utilizó el circuito propuesto por el profesor que se puede ver en la figura \ref{crkt:rectificador}.

\begin{figure}[!ht]
  \centering
  \begin{tikzpicture}[transform shape]
	% Paths, nodes and wires:
	\node[op amp, yscale=-1] at (0.083, 2.002){};
	\node[op amp, yscale=-1] at (-0.06, -2.49){};
	\draw (-6.5, 2.492) to[empty Schottky diode, l={$MBD101$}] (-4.75, 2.492);
	\draw (-4.5, -2.98) to[empty Schottky diode, l_={$MBD101$}] (-6.5, -2.98);
	\draw (-3.5, 2.5) to[american resistor, l={$100k\Omega$}] (-3.5, 1);
	\node[ground] at (-3.5, 1){};
	\draw (-1.107, 1.512) |- (1.5, 0.5) -| (1.5, 2) |- (1.273, 2.002);
	\draw (3, 2) to[american resistor, l_={$56k\Omega$}] (3, 0.25);
	\draw (3, -2.5) to[american resistor, l={$56k\Omega$}] (3, -0.75);
	\draw (3, -2.49) -| (1.13, -2.48);
	\draw (3, -0.75) -| (3, 0.25);
	\draw (3, -2.5) |- (2.75, -4.25) |- (-0.25, -4.25);
	\draw (-0.25, -4.25) to[american resistor, l_={$22k\Omega$}] (-2, -4.25);
	\draw (-2, -4.23) -| (-2, -2.98) -- (-1.25, -2.98);
	\draw (-4.25, -3) to[american resistor, l={$100k\Omega$}] (-4.25, -5);
	\node[ground] at (-4.25, -5){};
	\draw (-3, -2) |- (-1.25, -2);
	\node[ground, rotate=-90] at (-3.25, -2){};
	\draw (-3, -2) -- (-3.25, -2);
	\draw (-4.75, 2.492) -| (-1.107, 2.484);
	\draw (-6.5, -2.98) -- (-6.5, 2.492);
	\draw (1.5, 2.002) -| (3, 2);
	\draw (-4.25, -2.98) to[american resistor, l={$22k\Omega$}] (-2, -2.98);
	\draw (-4.25, -3) -| (-4.25, -2.98);
	\draw (-4.5, -2.98) -- (-4.25, -2.98);
	\node[ocirc](N1) at (-7, -0.25){} node[anchor=east] at ([xshift=-0.17cm]N1.east){$V_{2nd-derivated}$};
	\node[ocirc](N2) at (3.75, -0.25){} node[anchor=west] at ([xshift=0.2cm]N2.west){$V_{rectified}$};
	\draw (-7, -0.25) -- (-6.5, -0.25);
	\draw (3, -0.25) -- (3.75, -0.25);
\end{tikzpicture}
  \caption{circuito amplificador operacional en configuración rectificadora.}
  \label{crkt:rectificador}
\end{figure}

La simulación de esta etapa dio como resultado la señal que se puede ver en la figura \ref{plot:rectified_pre}.

\begin{figure}[!ht]
  \centering
  \pgfplotsset{
    table/col sep=semicolon,
    /pgf/number format/read comma as period
}

\begin{tikzpicture}
    \begin{axis}[
        width=.9\textwidth,
        height=5cm, % Reducido un poco más para seguridad
        xlabel={Tiempo [ms]},
        grid=both,
        xmin=0.00905, xmax=0.00952,
        scaled x ticks=false,
        xticklabel={\pgfmathparse{\tick*1000}\pgfmathprintnumber{\pgfmathresult}},
        x tick label style={/pgf/number format/fixed, /pgf/number format/precision=2},
        ytick = {0,0.2,0.4,\pgfkeysvalueof{/pgfplots/ymax}},
        %extra y ticks={\pgfkeysvalueof{/pgfplots/ymax}},
        ylabel={Amplitud [V]}
    ]
        \addplot[blue, thick] table [x=time, y={V(/rectified_pre)}] {sources/plots/sim.csv};
    \end{axis}
\end{tikzpicture}


  \caption{señal doble derivada rectificada.}
  \label{plot:rectified_pre}
\end{figure}

Nuevamente, esta señal de salida se encuentra fuera del rango solicitado, por lo que podríamos usar las ecuaciones
previamente utilizadas para el amplificador en configuración diferenciador, pero como no hay necesidad de cambiar el
offset ni cambiar la fase, basta simplemente con utilizar el amplificador no inversor para llegar al potencial pico
deseado.

\begin{figure}[!ht]
  \centering
  \begin{subfigure}[c]{0.6\textwidth}
    \begin{tikzpicture}[transform shape]
	% Paths, nodes and wires:
	\node[op amp, yscale=-1] at (-0, 0){};
	\draw (0.667, -1.752) to[american resistor, l={$R_2$}] (-1.083, -1.752);
	\draw (-3.583, -0.49) to[american resistor, l={$R_1$}] (-1.833, -0.49);
	\draw (1.19, 0) -| (2.917, -0.002);
	\node[ocirc](N1) at (2.917, 0){} node[anchor=west] at ([xshift=0.2cm]N1.west){$v_o$};
	\draw (-1.19, -0.49) -- (-1.833, -0.49);
	\draw (-1.583, -0.502) -- (-1.583, -1.752) -- (-1.083, -1.752);
	\draw (0.667, -1.752) -| (1.417, -0.002);
	\draw (-1.19, 0.49) -| (-3.833, 0.498);
	\draw (-3.583, -0.502) -- (-3.583, -2.252);
	\node[ground] at (-3.583, -2.252){};
	\node[ground] at (-3.583, -2.252){};
	\node[ocirc](N2) at (-3.833, 0.498){} node[anchor=east] at ([xshift=-0.2cm]N2.east){$v_i$};
\end{tikzpicture}
    \caption{circuito.}
    \label{crkt:amplificador.no.inversor}
  \end{subfigure}
  \hfill
  \begin{subfigure}[c]{0.3\textwidth}
    \begin{equation}
      v_o = v_i \frac{R_2}{R_1}
      \label{eq:amplificador.no.inversor.vo}
    \end{equation}
    \caption{ecuación.}
  \end{subfigure}
  \caption{amplificador operacional en configuración no inversora.}
\end{figure}

Como el pico de tensión que hay a la salida del rectificador es de $0.63 \si{\volt}$ y el potencial objetivo es de $1
\si{\volt}$ basta con encontrar la relación de ganancia:

\begin{align*}
  A_v &= \frac{v_o}{v_i} \\
  A_v &= \frac{1 \si{\volt}}{0.63 \si{\volt}} \\
  A_v &= 1.58
\end{align*}

El juego de resistencias que cumplen esta relación son:

\begin{equation*}
  R_1 = 56 \si{\kilo \ohm} \quad R_2 = 33 \si{\kilo \ohm}
\end{equation*}

Colocando la señal de salida del integrador en la entrada inversora y la tensión de referencia negativa en la entrada no
inversora, la simulación dio los resultados que se pueden ver en la figura \ref{plot:rectified}.

\begin{figure}[!ht]
  \centering
  \input{figures/plot_rectified.tex}
  \caption{señal rectificada con amplitud correcta.}
  \label{plot:rectified}
\end{figure}

\section{Señales Resultantes}
El esquemático completo está anexado al final del informe.

\begin{figure}[!ht]
    \centering
    \usepgfplotslibrary{groupplots}

\pgfplotsset{
    table/col sep=semicolon,
    /pgf/number format/read comma as period
}

\begin{tikzpicture}
    \begin{groupplot}[
        group style={
            group size=1 by 6,
            vertical sep=1cm, % Reducido para ahorrar espacio
            x descriptions at=edge bottom,
        },
        width=.9\textwidth,
        height=3.5cm, % Reducido un poco más para seguridad
        xlabel={Tiempo [ms]},
        grid=both,
        xmin=0.00905, xmax=0.00952,
        scaled x ticks=false,
        xticklabel={\pgfmathparse{\tick*1000}\pgfmathprintnumber{\pgfmathresult}},
        x tick label style={/pgf/number format/fixed, /pgf/number format/precision=2},
    ]
        % 1. BUFFER
        \nextgroupplot[ylabel={Buffer [V]}, title={1. Etapa de Buffer}]
        \addplot[blue, thick] table [x=time, y={V(/vi_buffered)}] {sources/plots/sim.csv};
        
        % 2. LEVEL SHIFT
        \nextgroupplot[ylabel={Offset [V]}, title={2. Etapa de Offset}]
        \addplot[blue!70!black, thick] table [x=time, y={V(/vi_offseted)}] {sources/plots/sim.csv};
        
        % 3. INTEGRATOR
        \nextgroupplot[ylabel={Integrador [V]}, title={3. Etapa de Integración}]
        \addplot[red, thick] table [x=time, y={V(/integrated)}] {sources/plots/sim.csv};

        % 4. DERIVATOR 1
        \nextgroupplot[ylabel={Deriv. 1 [V]}, title={4. Primera Derivación}]
        \addplot[magenta, thick] table [x=time, y={V(/derivated)}] {sources/plots/sim.csv};

        % 5. DERIVATOR 2
        \nextgroupplot[ylabel={Deriv. 2 [V]}, title={5. Segunda Derivación}]
        \addplot[cyan, thick] table [x=time, y={V(/derivated_twice)}] {sources/plots/sim.csv};

        % 6. RECTIFIER
        \nextgroupplot[ylabel={Rectif. [V]}, title={6. Etapa de Rectificación}]
        \addplot[green!60!black, thick] table [x=time, y={V(/rectified)}] {sources/plots/sim.csv};
        
    \end{groupplot}
\end{tikzpicture}

    \caption{Resultados de la simulación transitoria para todas las etapas del sistema.}
    \label{fig:plot_simulacion}
\end{figure}

